\chapter{MTHFR Gene Testing}

\section*{What Is This Test?}
MTHFR gene testing detects polymorphisms in the methylenetetrahydrofolate reductase (MTHFR) gene located on chromosome 1p36.3. The MTHFR enzyme catalyzes the conversion of 5,10-methylenetetrahydrofolate to 5-methyltetrahydrofolate, the circulating form of folate and the methyl donor for the remethylation of homocysteine to methionine. The two most commonly tested MTHFR polymorphisms are C677T (c.665C>T, rs1801133, alanine to valine substitution) and A1298C (c.1286A>C, rs1801131, glutamate to alanine substitution). The C677T homozygous genotype (TT, thermolabile variant) reduces MTHFR enzyme activity to approximately 30\% of normal, resulting in mildly elevated plasma homocysteine levels and mildly reduced folate levels. The A1298C homozygous genotype (CC) reduces enzyme activity to approximately 60\% of normal. Compound heterozygosity (C677T/A1298C) reduces activity to approximately 40--50\%. MTHFR genetic testing is performed by PCR and allele-specific probe detection (TaqMan, melt curve analysis, or direct Sanger sequencing). IMPORTANT NOTE: the American College of Medical Genetics and Genomics (ACMG), the American College of Obstetricians and Gynecologists (ACOG), and the American Heart Association do NOT recommend routine MTHFR genetic testing because the association between MTHFR polymorphisms and clinical outcomes (venous thromboembolism, recurrent pregnancy loss, neural tube defects, cardiovascular disease) is weak and inconsistent. Plasma homocysteine measurement is the clinically relevant test when homocysteine-related vascular or obstetric complications are suspected; MTHFR genotyping does NOT add clinically useful information beyond the homocysteine level alone.

\section*{Alternative Names}
MTHFR Mutation Analysis; MTHFR C677T and A1298C Genotyping; Methylenetetrahydrofolate Reductase Polymorphism; Thermolabile MTHFR; MTHFR DNA Test; MTHFR Gene Mutation

\section*{Principle}
MTHFR genotyping is performed on genomic DNA extracted from peripheral blood. Real-time PCR with allele-specific fluorescent probes (TaqMan assay) detects the presence of the wild-type (C or A) and variant (T or C) alleles at the C677T and A1298C loci. The fluorescence signal pattern distinguishes homozygous wild-type, heterozygous, and homozygous variant genotypes. Melting curve analysis on the LightCycler or similar real-time PCR platform is another common method: post-amplification melting curve analysis of a fluorescent hybridization probe identifies the variant allele by a characteristic lower melting temperature. Sanger sequencing of the relevant exons provides definitive confirmation. Results are reported as: C677T: C/C (wild-type, normal enzyme activity), C/T (heterozygous, approximately 60--70\% activity), T/T (homozygous variant, approximately 30\% activity); A1298C: A/A (wild-type), A/C (heterozygous), C/C (homozygous variant, approximately 60\% activity). Compound heterozygous: C677T C/T + A1298C A/C.

\section*{History}
MTHFR enzyme deficiency was first described in 1972 as a rare inborn error of metabolism causing severe hyperhomocysteinemia, homocystinuria, and neurologic disease. In 1987--1988, the thermolabile MTHFR variant (C677T) was identified by Kang et al. and shown to be a common polymorphism (approximately 10--15\% homozygous in Caucasian populations) associated with mild hyperhomocysteinemia. In the 1990s, numerous association studies reported links between MTHFR C677T and increased risk of venous thromboembolism, neural tube defects, recurrent pregnancy loss, coronary artery disease, and stroke. However, larger, better-controlled studies and meta-analyses published in the 2000s substantially weakened or refuted many of these associations, particularly after accounting for folate status. The 1998 mandatory folic acid fortification of cereal grains in the United States significantly reduced the population prevalence of neural tube defects and eliminated MTHFR-related hyperhomocysteinemia as a clinical concern in folate-replete populations. In 2013, the ACMG published an evidence-based statement concluding that MTHFR genotyping has no clinical utility for the evaluation of thrombophilia, recurrent pregnancy loss, or cardiovascular risk. In 2018, the American Heart Association similarly recommended against MTHFR testing for cardiovascular risk assessment.

\section*{Why This Test Is Ordered}
MTHFR testing was historically ordered for evaluation of unexplained venous thromboembolism (VTE), recurrent pregnancy loss, preeclampsia, placental abruption, fetal neural tube defects, and hyperhomocysteinemia. HOWEVER, current guidelines DO NOT recommend MTHFR genetic testing for any of these indications. Plasma homocysteine measurement is the appropriate test when hyperhomocysteinemia is suspected. In the very rare situation of severe hyperhomocysteinemia ($>$100 mcmol/L) in a patient with homocystinuria, MTHFR enzyme deficiency can be considered, but this is a severe inborn error of metabolism distinct from the common C677T polymorphism. The current, evidence-based indications for NOT ordering MTHFR testing are: venous thromboembolism workup (ACMG 2013), recurrent pregnancy loss (ACOG 2020), cardiovascular disease risk assessment (AHA 2018), and isolated mild hyperhomocysteinemia (which is managed by folate supplementation regardless of MTHFR genotype).

\section*{Primary vs Secondary Test}
MTHFR genetic testing is currently NOT recommended as a primary or secondary test by any major medical society for any common clinical indication. Plasma homocysteine, when clinically indicated, is the appropriate laboratory test.

\section*{How This Test Is Performed}
For most patients, a health care professional will take a blood sample from a vein in your arm using a small needle. After the needle is inserted, a small amount of blood will be collected into a test tube or vial. You may feel a little sting when the needle goes in or out. This usually takes less than five minutes. DNA is extracted from leukocytes and genotyped by PCR.

\section*{What Preparation Is Needed}
No special preparation is required. Fasting is not necessary. If plasma homocysteine is being measured concurrently, an 8--12 hour fast may be recommended (homocysteine levels rise after a high-protein meal).

\section*{Contraindications and Precautions}
The most important clinical precaution is that MTHFR genotyping should NOT be ordered for routine evaluation of thrombophilia, recurrent pregnancy loss, or cardiovascular disease, as the test has no proven clinical utility for these indications. The C677T variant is very common: approximately 40--50\% of the population is heterozygous (C/T), and 8--15\% is homozygous (T/T) in Caucasian populations (prevalence varies by ethnicity: approximately 10--12\% T/T in Caucasians, 2--5\% in African Americans, $\sim$20\% in Hispanics). Because the homozygous variant is so common, many patients are found to be homozygous T/T, leading to unnecessary anxiety and non-evidence-based interventions such as aspirin or anticoagulation therapy (not recommended). MTHFR polymorphisms are NOT thrombophilias and do NOT require anticoagulation or antiplatelet therapy. The association between MTHFR C677T and venous thromboembolism was confounded by mild homocysteine elevation in folate-deplete populations and is largely eliminated in folate-replete populations (post-1998 folic acid fortification). MTHFR testing should not be performed in patients with known severe hyperhomocysteinemia (suspected inborn error), for whom direct homocysteine measurement and enzyme activity assays are the appropriate tests.

\section*{Risks}
Minimal physical risk from venipuncture. The greater risk is iatrogenic harm from misinterpretation of the test result: unnecessary medical interventions, inappropriate prescribing of anticoagulation or antiplatelet therapy (exposing the patient to bleeding risk), and psychological distress from a perceived ``genetic disease'' that, in fact, is a common benign normal variant.

\section*{What Happens After the Test}
Results are typically available in 3--7 days. A homozygous or compound heterozygous result should NOT trigger any specific medical intervention. The appropriate clinical response is reassurance. If plasma homocysteine is concurrently elevated (>15 mcmol/L), folic acid supplementation (1 mg/day), vitamin B12, and vitamin B6 may be recommended regardless of MTHFR genotype, as the treatment of hyperhomocysteinemia is the same regardless of MTHFR status. Patients should be counseled that the variant is a common normal variant and does not confer an increased risk of thrombosis, adverse pregnancy outcomes, or cardiovascular disease.

\section*{Understanding Results}

\subsection*{C677T Homozygous Wild-Type (C/C)}
Normal MTHFR enzyme activity (approximately 100\%). Most common genotype (40--50\% of the population). No clinical implications.

\subsection*{C677T Heterozygous (C/T)}
Mildly reduced MTHFR enzyme activity (approximately 60--70\%). Found in 40--50\% of the population. Typically normal plasma homocysteine in folate-replete individuals. No proven adverse clinical effects. No intervention is required.

\subsection*{C677T Homozygous Variant (T/T)}
Moderately reduced MTHFR enzyme activity (approximately 30\% at 37 degrees C). Found in 8--15\% of the population. Increased plasma homocysteine (typically mild, 12--20 mcmol/L) restricted to individuals with insufficient folate intake. In the era of universal folic acid fortification of grain products (USA, 1998), the clinical impact of this genotype on homocysteine levels has been substantially attenuated. No independent association with VTE, recurrent pregnancy loss, or cardiovascular disease in folate-replete populations. No indication for anticoagulation or antiplatelet therapy. No alteration in pregnancy management. ACOG specifically states that MTHFR C677T TT is NOT an indication for thromboprophylaxis during pregnancy.

\subsection*{A1298C Homozygous Variant (C/C)}
Mildly reduced MTHFR enzyme activity (approximately 60\%). No consistently demonstrated association with hyperhomocysteinemia or any adverse clinical outcome. No clinical implications.

\section*{Normal Results}
MTHFR C677T: C/C genotype (wild-type), normal MTHFR enzyme activity (approximately 100\%). MTHFR A1298C: A/A genotype (wild-type), normal MTHFR enzyme activity. Normal plasma homocysteine: 5--15 mcmol/L in adults. MTHFR testing is NOT recommended as a routine test by ACMG, ACOG, or the AHA. Abnormal results (homozygous or compound heterozygous) are common normal variants and do not indicate disease.

\section*{Follow-up Tests}
Plasma homocysteine: the clinically relevant test. Fasting plasma homocysteine >15 mcmol/L indicates hyperhomocysteinemia and warrants investigation for vitamin deficiencies (folate, vitamin B12, vitamin B6), renal insufficiency, hypothyroidism, or medications (methotrexate, phenytoin, nitrous oxide). Serum folate and vitamin B12 levels: to identify vitamin deficiency as a cause of hyperhomocysteinemia. Factor V Leiden and prothrombin G20210A mutation testing: if a true hereditary thrombophilia workup is indicated for recurrent VTE, NOT MTHFR. Antiphospholipid antibody testing: if recurrent pregnancy loss is being evaluated. Homocystinuria evaluation: plasma amino acid analysis in an infant or child with severe hyperhomocysteinemia >100 mcmol/L, ectopia lentis, developmental delay, or thromboembolism. MTHFR enzyme activity in cultured fibroblasts: if severe MTHFR deficiency (the rare inborn error of metabolism) is suspected.

\section*{References}
Hickey SE, Curry CJ, Toriello HV. ACMG Practice Guideline: lack of evidence for MTHFR polymorphism testing. Genet Med. 2013;15(2):153--156. Practice Bulletin No. 197: Inherited Thrombophilias in Pregnancy. American College of Obstetricians and Gynecologists. Obstet Gynecol. 2018;132(1):e18--e34. Greenland P, Alpert JS, Beller GA, et al. 2010 ACCF/AHA guideline for assessment of cardiovascular risk in asymptomatic adults. J Am Coll Cardiol. 2010;56(25):e50--e103. Clarke R, Bennett DA, Parish S, et al. Homocysteine and coronary heart disease: meta-analysis of MTHFR case-control studies, avoiding publication bias. PLoS Med. 2012;9(2):e1001177. Levin BL, Varga EA. MTHFR: addressing genetic counseling dilemmas using evidence-based literature. J Genet Couns. 2016;25(5):901--911.

\clearpage
\section*{Regulatory Status and Authorization Date}
Regulatory status applies to the specific assay, instrument, manufacturer, and intended use rather than to the test name alone. No single approval date can be assigned to this test category. For a commercial assay, verify the current product in the relevant regulator's database: the U.S. Food and Drug Administration (FDA) clearance, approval, or authorization; India's Central Drugs Standard Control Organisation (CDSCO) licence or permission; the European Union In Vitro Diagnostic Regulation (EU IVDR) CE certificate issued through the applicable conformity-assessment route; or the United Kingdom Medicines and Healthcare products Regulatory Agency (MHRA) registration and UKCA/CE status. Laboratory-developed tests may not have an individual FDA, CDSCO, EU IVDR, or MHRA approval date.
\clearpage
